Across many markets, firms buy data directly from consumers---with attractive rewards for data-sharing and readily available sensor technologies---and keep what they collect proprietary. The data is then used to mitigate information asymmetries, strengthen competitive advantage, and extract rents from consumers. We obtain novel datasets on a voluntary monitoring program in the competitive U.S. auto insurance industry, which give us direct visibility into how proprietary monitoring data is collected and used beyond prior empirical analyses that have relied on simulations \parencites{Bordoff2008} and state-firm level aggregates \parencites{Hubbard2000, Reimers2018, porrini2020black}. We develop an empirical framework that embeds the firm's data elicitation strategy, and consumers' data-sharing choices, inside the broader product (insurance) market pricing and choice problems. This allows us to account for the value of monitoring data to the firm based on its subsequent use in price discrimination, risk-rating, competitive cream-skimming, and risk reduction.

We first separately identify moral hazard from selection. This extends the literature that has relied on correlations between plan choice and claims, such as \textcite{Chiappori2000, Cohen2007}. \textcite{abbring2008better, jeziorski2019skimming} examine moral hazard through dynamic incentive variation in experience-rated pricing, where surcharges for further accidents rise with each accident. Moral hazard thus implies declining claim rates after an accident. However, recent evidence suggests that drivers become inherently more cautious following ``near misses,'' even when no incentive changes occur\parencite{shum2022time}.

Our work builds on a structural literature measuring the extent and welfare implications of information asymmetries in insurance markets \parencites{einav2010estimating}. Like \textcite{Cohen2007}, \textcite{Barseghyan2013b}, and \textcites{jeziorski2019skimming}, we use consumer choices and subsequent claims data from an auto insurer to identify consumers' accident risk and risk preferences. Our estimates highlight a similar information asymmetry between the insurer and its consumers, which drives advantageous selection into monitoring. Beyond modeling selection, the monitoring program allows both the firm's information set and that of its competitors (in our counterfactuals) to be determined endogenously in equilibrium with insurance prices and quantities. To do this, we expand the canonical framework in two main ways.

First, the revelation of risk information takes time. When consumers decide whether or not to opt in to monitoring, they anticipate the effect that the monitoring score they receive will have on their renewal price. Our choice model reflects the dynamic nature of monitoring decisions, as well as uncertainty about future premiums that consumers face given their risk types and risk preferences. This allows us to predict how changes in opt-in discounts and renewal pricing alter selection into monitoring (i.e. how many and what kinds of consumers choose to opt in), and consequently, the amount of information that is revealed.

Second, monitoring data is proprietary, and firm pricing must account not only for selection into the program but also for selection into the firm. Accounting for competitor pricing is thus crucial to understanding consumers' outside options and the firm's pricing incentives. Compared to existing studies, we introduce new individual-level competitor pricing data based on price filings; we also expand the canonical model to admit consumer choice along both the intensive and extensive margins. Similar work include \textcite{Honka2012} that uses survey data and \textcite{cosconati2025competing}, which uses administrative data.

We further contribute to the empirical literature on the role of data in industrial organization.
First, monitoring mirrors more traditional modes of hard information disclosure such as restaurant health score cards \parencites{jin2003effect} and used-car photos \parencites{lewis2011asymmetric}. Our study differs conceptually in that consumers have control over information disclosure, although their disclosure decisions are influenced by firm pricing. Nonetheless, our findings highlight similar sources of efficiency benefits from richer information as well as the signaling incentives and frictions related to its disclosure.
Second, our final counterfactual considers a regulation that would require public disclosure of monitoring data, which connects to recent policy debates regarding data portability and algorithmic transparency \parencites{jin2021big}. We find that the loss of proprietary control over monitoring data not only heightens the reclassification risk associated with monitoring but also erodes the firm's ability to recoup its upfront ``investment,'' consisting primarily of the cost of opt-in discounts used to encourage participation. Imposing a minimum discount floor would exacerbate these effects, further reducing consumer welfare.
This underscores the trade-off between curbing markups and protecting the firm's incentive to produce data in the first place \parencites{Posner1978, Hermalin2006}.

Finally, our findings provide empirical support for a recent theoretical literature on voluntary disclosure and personalized pricing.
Monitoring data is a form of \textit{hard} (verifiable) information that cannot easily be falsified. Although voluntary disclosure of such information by some consumers may lead a monopolist firm \parencite{pram2020disclosure}, or one facing horizontal competition \parencite{alv2022voluntarydisclosure}, to infer that others have lower types, consumers can still obtain a Pareto improvement, corresponding to more coverage purchased and hence more risk-sharing in our setting, if there are gains from trade with consumer types that wouldn't be served without the additional information. Consumer control over information sharing, realized through the ex-ante opt-in structure in our setting, mitigates consumer harm from price discrimination. As in \textcites{ichihashi2020online, montes2019value}, we argue that consumer control may thus obviate the need for ex-post regulations on exclusive data ownership.

Since our data period (2012--2016), the monitoring landscape has changed substantially. Most major U.S. auto insurers now offer smartphone-based telematics programs that run in the background and automatically evaluate driving trips, lowering the marginal cost of data collection relative to the OBD-device era.\footnote{See Section \ref{sec:bg_data} for a description of the technological transition. As of 2024, industry estimates suggest that 12--22\% of U.S. drivers are enrolled in a telematics program (J.D.\ Power, 2025; AutoInsurance.com, 2024; \cite{ben2023privacy}).} Yet adoption remains limited. Several forces documented in our paper help explain this. The monitoring disutility that we estimate captures privacy loss, increased driving effort, and additional decision-making costs; these frictions remain a participation barrier \parencites{athey2017digital}{ben2023privacy}. Regulatory constraints continue to bind: California effectively prohibits behavioral telematics pricing, while states such as New York restrict monitoring adjustments to discounts only \parencite{ben2023privacy}. Monitoring data also remains proprietary and non-portable across insurers, unlike claims records, which are tracked in centralized industry databases for up to seven years. At the same time, the broader evidence on behavioral monitoring in transportation, from onboard computers in trucking \parencite{Hubbard2000} to recent work using continuous telematics data \parencites{soleymanian2019sensor}{jin2023re}, reinforces our finding that even imperfect monitoring signals can generate meaningful safety improvements. The fundamental trade-offs we identify between moral hazard reduction, adverse selection, reclassification risk, and participation frictions are inherent to any monitoring program regardless of the underlying technology. Our finding that monitoring generates positive total surplus is reinforced by the lower data collection costs and broader program availability observed in today's market.
